\begin{abstract}
\textbf{Background.} The dominant sequence transduction models are based on complex recurrent or convolutional neural networks with encoder-decoder architectures. While attention mechanisms have been integrated, the sequential nature of recurrence limits parallelization during training, becoming increasingly problematic for longer sequences.
\textbf{Methods.} We propose the Transformer, a model architecture eschewing recurrence and instead relying entirely on an attention mechanism to draw global dependencies between input and output. The model employs stacked self-attention and point-wise, fully connected layers for both encoder and decoder, with multi-head attention allowing joint attention to information from different representation subspaces.
\textbf{Results.} Experiments on the WMT 2014 English-to-German translation task achieve 28.4 BLEU, improving over existing best results by over 2 BLEU. On the WMT 2014 English-to-French translation task, the model establishes a new state-of-the-art BLEU score of 41.0 with significantly reduced training cost compared to prior models.
\textbf{Conclusion.} The Transformer architecture demonstrates that attention mechanisms alone can effectively model sequence transduction, enabling significantly faster training than recurrent or convolutional layer-based approaches. This work opens new possibilities for parallelizable sequence modeling and may inspire future applications beyond translation.
\\
\textbf{Keywords:} Transformer, Self-Attention, Machine Translation, Multi-Head Attention, Sequence Modeling
\end{abstract}

\selectlanguage{chinese}
\begin{abstract}
\textbf{Background.} 主流的序列转换模型基于复杂的循环或卷积神经网络，采用编码器-解码器架构。虽然已集成注意力机制，但循环的序列特性限制了训练期间的并行化能力，对较长序列而言问题愈发突出。
\textbf{Methods.} 我们提出了Transformer模型架构，摒弃循环结构，完全依赖注意力机制来建模输入和输出之间的全局依赖关系。该模型在编码器和解码器中均使用堆叠的自注意力和逐点全连接层，通过多头注意力机制实现对不同表示子空间信息的联合关注。
\textbf{Results.} 在WMT 2014英德翻译任务上达到28.4 BLEU，比现有最佳结果提升超过2 BLEU。在WMT 2014英法翻译任务上，模型以显著降低的训练成本，取得41.0 BLEU的全新单模型最优成绩。
\textbf{Conclusion.} Transformer架构证明，仅靠注意力机制就能有效建模序列转换，相比基于循环层或卷积层的方法可实现显著更快的训练。这项工作为可并行化的序列建模开辟了新的可能，并可能启发翻译之外的更多应用。
\\
\textbf{Keywords:} Transformer, Self-Attention, Machine Translation, Multi-Head Attention, Sequence Modeling
\end{abstract}

% BibTeX Entry
@article{attentionis,
  title = {Attention Is All You Need},
  abstract = {The dominant sequence transduction models are based on complex recurrent or convolutional neural networks with encoder-decoder architectures. While attention mechanisms have been integrated, the sequential nature of recurrence limits parallelization during training, becoming increasingly problematic for longer sequences. We propose the Transformer, a model architecture eschewing recurrence and instead relying entirely on an attention mechanism to draw global dependencies between input and output. The model employs stacked self-attention and point-wise, fully connected layers for both encoder and decoder, with multi-head attention allowing joint attention to information from different representation subspaces. Experiments on the WMT 2014 English-to-German translation task achieve 28.4 BLEU, improving over existing best results by over 2 BLEU. On the WMT 2014 English-to-French translation task, the model establishes a new state-of-the-art BLEU score of 41.0 with significantly reduced training cost compared to prior models. The Transformer architecture demonstrates that attention mechanisms alone can effectively model sequence transduction, enabling significantly faster training than recurrent or convolutional layer-based approaches. This work opens new possibilities for parallelizable sequence modeling and may inspire future applications beyond translation.},
  keywords = {Transformer; Self-Attention; Machine Translation; Multi-Head Attention; Sequence Modeling},
  note = {Methods: Transformer Architecture; Multi-Head Self-Attention; Position-Wise Feed-Forward Networks; Residual Connections},
}